\justifying
\NSFCBodyText

\subsubsubsection{支撑创新一：三时钟观测下的条件辨识}

本项目分离事件发生、设备采样和决策端到达过程，使“未采”与“未达”保留不同的风险含义。可观测锚点和支持条件成立时估计事件风险；条件偏离时用敏感性区间与拒判暴露不可辨识性，并以同等时间信息强基线和条件破坏实验检验。

\subsubsubsection{主创新：集合条件证据价值与安全资源决策}

项目把候选证据加入现有集合并在截止期前交付后的风险降低作为统一决策量，使互补、冗余和迟到失效进入同一价值估计。高可靠母流和共同随机数重放提供反事实核验，真实日志只在行为支持域内使用；策略仅在关键原始证据保留、最低告警、压缩保真和幂等续传构成的安全动作集中优化。该机制区别于按单包置信度、新鲜度或大小排序，是项目的主创新。

\subsubsubsection{支撑创新二：漂移校准与条件性事件记忆}

仓储部分先以本地温湿度、烟雾、门禁和视频风险的漂移校准、选择性预测与分级响应构成独立主问题。货物级配对和统计功效满足后，再检验冲击/倾斜或非法开启摘要的事件特异衰减及组外增益，并用摘要置换、时间错配和支持域外拒判排除普通历史拼接造成的伪增益。
